%---------------------------------------------------------------------------
% Cold-outreach letter template — same brand as cover letter
% Per Eugene Vinitsky's three principles + v3 outreach research:
%   1. Demonstrate familiarity (specific paper, specific section)
%   2. Make a SPECIFIC ask (15 min, a pointer, a question)
%   3. Provide credentials if junior (ONE link, in signature)
% Body should be UNDER 80 words.
%---------------------------------------------------------------------------
\documentclass[11pt,letterpaper]{article}

\usepackage{cmap}
\usepackage[T1]{fontenc}
\usepackage[utf8]{inputenc}
\usepackage[margin=1.0in,top=0.7in,bottom=0.85in]{geometry}

\usepackage[scale=0.97]{plex-serif}
\usepackage[scale=0.91]{plex-sans}
\usepackage[scale=0.89]{plex-mono}
\usepackage[protrusion=true,expansion=true,tracking=true,kerning=true,final]{microtype}
\AtBeginDocument{\figureversion{osf}}
\frenchspacing

\linespread{1.20}
\setlength{\parindent}{0pt}

\usepackage{xcolor}
\definecolor{accent}{HTML}{0A2540}
\definecolor{accentlight}{HTML}{2563A8}
\definecolor{meta}{HTML}{6B6B6B}
\definecolor{rule}{HTML}{C8C8C8}
\color{black!90}

\usepackage[hidelinks, colorlinks=true, urlcolor=accentlight, linkcolor=accentlight]{hyperref}
\usepackage{xurl}
\usepackage{parskip}
\setlength{\parskip}{0.6\baselineskip}

\newcommand{\bdash}{\,\textemdash\,}

\begin{document}

%---------------------------------------------------------------------------
% This is a TEMPLATE. The actual cold email is the body text below; the
% PDF version exists only for cases where the recipient explicitly asked
% for an attachment OR you want a print-quality copy of the same content.
%---------------------------------------------------------------------------

\textbf{Subject:} Question on \S[3.2] of ``[Paper Title]''

\vspace{1em}

Hi [First Name],

\vspace{0.3em}

Your [paper / talk] on [specific topic] showed [one concrete fact: e.g.,
``the modality gap collapses when you co-train with the masked
objective'']. I'm working on [one-line: e.g., ``tabular foundation
models with image-column inputs''] and ran into [specific connection:
e.g., ``the same gap pattern but only on heterogeneous schemas''].

\vspace{0.3em}

Could I send you a 1-page write-up of what I found and ask one
question? Happy to do it async \bdash{} no call needed unless you'd prefer.

\vspace{0.3em}

Dennis Loevlie

\vspace{0.5em}

{\sffamily\footnotesize\color{meta}%
ELLIS PhD candidate \bdash{} CWI (TRLLab) \& UvA (AMLab) \bdash{} \\
\href{https://dennisloevlie.com}{dennisloevlie.com} \bdash{}
CV: \href{https://loevlie.github.io/cv/loevlie-cv-latest.pdf}{loevlie.github.io/cv} \bdash{}
[arxiv link OR github repo of most relevant work]
}

\vspace{2em}
{\color{rule}\rule{\linewidth}{0.4pt}}
\vspace{1em}

{\sffamily\footnotesize\color{meta}%
\textit{Variants from the v3 outreach research:}%
\begin{itemize}
\item \textbf{Collaboration ask:} replace closing with ``I have [data / compute / a complementary result]; would a 4-week collaboration on [specific output, e.g., `a workshop paper at TRL'] interest you?''
\item \textbf{Referral / informational chat:} ``I'm applying for [program / position] this cycle. If you'd be open to a 15-minute call on [date range], I'd value your read on whether my [project X] would land at [team Y].''
\item \textbf{Post-talk follow-up:} ``Thanks for the talk at [venue] on [date] \bdash{} your point on [specific moment] connected to [your work]. [One question]. If you have 15 minutes in the next two weeks, I'd love to follow up.''
\end{itemize}

\textit{Anti-patterns to avoid:} ``I hope this email finds you well'' \bdash{} ``I've been a huge fan of your work'' \bdash{} ``Could I pick your brain'' \bdash{} multiple attachments on first email \bdash{} more than two follow-ups.

\textit{Subject line patterns that get opened:}\\
``Question on \S3.2 of [Paper]'' \bdash{}
``Reproduced your [X] result \bdash{} small finding'' \bdash{}
``After your NeurIPS talk: one question on [X]'' \bdash{}
``Collab idea: [your method] on [their benchmark]''
}

\bdash{}

\renewcommand{\bdash}{\,\textemdash\,}
\end{document}
